\documentclass[11pt]{article}

\usepackage{bbm}
\usepackage{geometry}
\usepackage{graphicx}
\usepackage{amsmath}
\usepackage{float}
\usepackage{listings}
\usepackage{indentfirst}
\usepackage{braket}
\usepackage{authblk}
\setlength{\parindent}{2em}
\geometry{left=2.5cm,right=2.5cm,top=2.0cm,bottom=2.0cm}
\usepackage{hyperref}
% About the author
\renewcommand*{\Authand}{, }
\renewcommand\Affilfont{\itshape\small}

\date{}

\title{Mathematical Model: From Cell Membrane to Simple Neural Network}
\author[1]{Zeyuan Ye}
\author[2]{Sebastian K{\"o}nig}
\setcounter{Maxaffil}{0}

\begin{document}
  \maketitle

\begin{abstract}

  Researches study our neural system in different levels: from
  conponents of a single cell, such as membrane, to the connection of two cells
  (synapsis), and further to the complex network
  of the whole brain. This summary trying to find a straight pathway of the
  modeling from sub-cellcular level to simple neural network. Please refer
  \cite{mitcourse} for more detail
  in.

\noindent \textbf{Keywords:} lorentz integral transformation; pionless EFT;
  bound-state problem\\ %(5 at most)
\end{abstract}

\section{Membrane}
Membrane has roughly two parts that are vital in the neural
activity: phospholipid bilayer and channels (FIG. \ref{structMemb}). phospholipid is permeterbale for water, oil
etc., yet impermmitable for ions. This attribution is similar to semi-permeable
membrane. Ions carry charges, if the charge
densities of the two sides of phospholipid bilayer (intra celluar space and
extra-cellcular space) are different, there will be a votage across cell
membrane (FIG. \ref{semiMemb}). In all, because of the existance of phospholipid bilayer, the cell has a
function similar to conductance
\begin{equation}
  c V = Q
  \label{capt}
\end{equation}
where $c$ is the conductance per unit area, $Q$ is the charge density
differences and $V$ is the votage accross the cell. The votage of the
extra-cellcular space is set to be zero, therefore $V$ here is also the votage
of intra-cellcular space, which is roughly $-70 \text{mV}$ to $ -50 \text{mV}$ as measured
by the experiment.

There are two types of channel: ion channels and metabolic channel. The effect
of metabolic channels are more indirect and complex, which is not the center
concern of this summary. Ion channel behaviours as the gate of the membrane. It
allows only certain types of ions to pass though the membrane, thus produce inward
or outward cuurent flow. For example,
Na-gated channel is only permitterble to Na$^+$, when the channel is open (FIG. \ref{nachannel}). One
good approximation of the channel is resistence, govern by the Ohm's law
\begin{equation}
  I = G_i (V - V_i)
  \label{olm}
\end{equation}
where $G_i$ is the conductance which can depends on the votage and time, and
$V_i$ is the equivibrilum votage (constant for a certain type of channel) results from
the balance of electric force and diffusion. Index $i$ here means
$ith$ type of ion channel. Positive $I$ means outward current.

\section{Single Cell}

Differentiate on both side of equation (\ref{capt}) and use equation ({\ref{olm}}), we
have an equation for single cell
\begin{equation}
  c \frac{dV}{dt} = - G_L (V - V_L) - I_e
  \label{ifmodel}
\end{equation}
where $I_e$ is the extra-induced current, such as the experimentist inject ions
to the cell (negtive $I_e$). Index $L$ is leakage channel, which is a
mathematically equvilent
channel of partial function of some types of channel (largely fall into K-selective channel). The rest of effects of channels can
be simplified as the following statement:
\begin{itemize}
\item{} If $V > V_{th}$, $V = 50 \text{mV}$, then quickly fall to $V_{res}$.
  \label{ifstate}
\end{itemize}
This effect is often introduced by Na-selective channels. Typical value of
$V_{res}$ is $-65 mV$. There are some choices for $V_{th}$, but values like $-30
mV$ is sufficient to fit a good result. When this statement is triggered, we say
the neuron fired, or it made a spike. This single-cell model also be named
Integrate and Firing Model. This model contains one resistence (leakage
channel), one conductance (membrane),
one external current source ($I_e$), and voltage source (FIG. \ref{ifmodel}). There are lots more sophiscated models, such as HH
model, which replace the statement with more realistic channels as
additional terms in equation (\ref{ifmodel}).

Equation \ref{ifmodel} can be collected with a more susinct form
\begin{equation}
  \tau \frac{dV}{dt} = -V + V_{\infty}
  \label{ifmodelS}
\end{equation}
where $\tau = c / G_L$ and $V_{\infty} = (V_L - I_e) / G_L$. If $G_L$ here is
a constant, solution of equation \ref{ifmodelS} is
\begin{equation}
  V = (V_L - V_{\infty}) e^{-t/\tau} + V_{\infty}
\end{equation}
The firing pattern is shown in FIG. \ref{fpattern}

Information is stored in the firing rate of 
neurons instead of the exact value of votage. Firing rate is defined as the
number of spiks in a unit time. The time spacing between two spikes can be
calculated from equation (\ref{ifmodelS}) by calculate the time that $V$ grows
from $V_{res}$ to $V_{th}$. The inverse of time spacing is firing rate, as show
in FIG. \ref{fireate}. In all, we could conclude the relation with
\begin{equation}
  v = F[I_e]
  \label{fireateeq}
\end{equation}
where $v$ is the firing rate, and $F$ is some function has the shape in FIG. \ref{fireate}

\section{Synapse}
Before moving to neural network, we need to understand how two neurons connect
each other by synapse. FIG. \ref{syn} shows the biological structure of the
synapse. When the pre-neuron is fired, the depolyzation will open the
Ca-channel, induces an inward flow of Ca$^{2+}$. This Ca$^{2+}$ will release the
transmitters in the presynapse. Transmitter in the synapse cleft then attarch to
particular recepters in the post-synapse, open ion-channel to allow Na$^+$ flow
into the post-synapse. This final lead to the depolyzation of the post-neuron.

The current flow into the post-neuron can be descripted by
\begin{equation}
  I_{syn} = - G_{syn}(t) (V - V_{syn}),
  \label{isyn}
\end{equation}
where $V_{syn}$ is about $0$ mV, which is the reversal potential of the
post-synapse. Conductance $G_{syn}(t)$ is shown in FIG. \ref{gsyn} If there's no
transmitters in the synapse cleft, ion channels in the post-synapse are closed,
$G_{syn}(t) = 0$ mV.
A spike of the presynapse will suddenly release a large
amount of transmitters, therefore largely increase
$G_{syn}(t)$. Then the transmitters will graduatly either diffuse away or be
decomposed by glials, which lead to the slow decay of the conductance.
The conductance induced by single spike is often approximated by the exponential
function
\begin{equation}
  K(t) = G_{max} e^{-t / \tau}.
\end{equation}
The whole conductance under a train of spikes is
\begin{equation}
  G(t) = \sum_iK(t) \delta_{t,t_i}
  \label{spikeraw}
\end{equation}
where there is a spike at time $t_i$. Equation (\ref{spikeraw}) can be
simplified with the spike train $S(t) = \sum_i\delta(t - t_i)$
\begin{equation}
  G(t) = \int{d\tau}K(\tau)S(t - \tau)
\end{equation}

\section{Simple Neural Network}
Consider a network with only two neurons attarched FIG. \ref{twonet}. The
working status of each neuron is solely stated by its firing rate. So we need to
findout the firing-rate relation between neurons. For the single neuron, the only source of external current is from its
pre-neuron, i.e. $v = F[I_{syn}]$ from equation (\ref{fireateeq}). Yet $I_{syn}$
followed by equation (\ref{isyn}) is still
too complex to be pratical. Hence we further throw away the votage term, and
treat $G_{syn}(t)$ as votage independent,
\begin{equation}
  I_{syn} = \int{d\tau}K(\tau)S(t - \tau)
\end{equation}
we could extract a factor $w$ so that the rest of $U(\tau) = K(\tau) / w$ can be normalized,
\begin{equation}
  \int{d\tau}U(\tau) = 1 
\end{equation}
The convolution of a narrow (correnspond to short time of non-zero $K(\tau)$)
window $U(\tau)$ with spike train $S(\tau)$ is exactly the firing rate of the
presynapse $u$. Overall, we have
\begin{equation}
  v = F[w u].
\end{equation}
This equation reveals the simple relation of the network. Larger $w$ means the
two neurons are more strongly connect, only a slow firing rate of the pre-neuron
is enough for triggering the post-neuron. This biologically could mean there are more
channels in the post-synapse. Different presynapse will attarch to different
sites of the postsynapse. Furthermore, the post-neual layer can also attarch to
each other as shown in FIG. . The overall effect is the summation of all their
current.
\begin{equation}
  \boldsymbol{v} = F[W\boldsymbol{u} + M \boldsymbol{v}],
\end{equation}
where $W, M$ are matrices, arrays are written as bold letters. This model is
called rate model, which is the fundation of the lots of deep learning
architechture.

\section{Conclusion}
Membrane of the neuron is a conbination of resistor and capacity. The
conductances are different under different membrane and more importantly,
different to different types of channels. This attribution allows the neuron to
fire, and soon go back to the rest potential. A spike
from the pre-neuron 
will triggered a bit depolyzation of the post-neuron. Once there is enough
depolyzation, the post-neuron will also make spikes.
 Despite the complicated procedure
for generating a spike, the exact voltage is unimportant. The real information is stored in the frequency of
the spikes, which can be described using a single number -- firing rate. By
carefully simplify the synapse conductance, one could obtain a relation between
the firing rate of neurons, thus have a neuron network. This neuron network
also lies in the basis current deep network algorithms.

\begin {figure}[H]
\centering 
\includegraphics[width=17cm]{test.png} 
\caption{results of problem2\_1}
\end {figure}

\begin{thebibliography}{100}
\bibitem{mitcourse}
\href{https://ocw.mit.edu/courses/brain-and-cognitive-sciences/9-40-introduction-to-neural-computation-spring-2018/}{MIT OpenCourseWare: Introduction to Neural Computation}
\end{thebibliography}
\end{document}